% Módulo 03b — Tutorial de estatística da calibração (didático, R773)
Números sem glossário não comunicam: como ler cada medida que aparece no Capítulo~6.
\section{Ler os números deste trabalho}
\textit{Arrependimento:} diferença entre a melhor utilidade possível e a do escolhido (0 = perfeito). \textit{Taxa de acerto:} fração de decisões exatamente ótimas. \textit{IC95:} faixa ``média $\pm$ margem''; ICs sem sobreposição indicam diferença além do ruído amostral, sob aproximação normal. \textit{Pareamento:} comparar estratégias nas mesmas 200 decisões elimina variação de dificuldade entre tarefas. \textit{Taxa de vitória:} fração de decisões em que a atenção empata ou supera a rival. \textit{Pearson $r$:} correlação linear entre $-1$ e $+1$; $r$ próximo de 0 indica ausência de relação linear (não de qualquer relação). \textit{Semente:} número que fixa todos os sorteios; mesmos números em reexecução provam determinismo. \textit{Oráculo:} verdade sintética oculta do roteador; seus pesos ($q_w$) são arbitrários --- por isso a sensibilidade em três mundos.
\section{Fundamentos com lastro}
Relato por intervalos segue a declaração sobre valores-\textit{p} \cite{wasserstein2016} e a nova estatística por estimação \cite{cumming2014}; reamostragens do tipo bootstrap \cite{efron1979} são a alternativa clássica quando a aproximação normal é suspeita. Escores de probabilidade bem calibrados obedecem a regras de escore próprias \cite{gneiting2007}.
\section{Exemplo numérico trabalhado}
A margem do arrependimento da atenção ($\pm0{,}0105$, $n=200$) implica desvio-padrão $s=0{,}0105\times\sqrt{200}/1{,}96\approx0{,}0758$; a diferença pareada contra pesos iguais ($+0{,}035\pm0{,}011$) mantém-se positiva em toda a faixa, logo robusta ao ruído amostral registrado.
\section{Derivação do intervalo (passo a passo)}
Sob aproximação normal, a média amostral $\bar{x}$ de $n$ observações com desvio $s$ varia em $s/\sqrt{n}$; o fator $1{,}96$ cobre $95\%$ da massa. Com $n=200$ e $s\approx0{,}0758$: $1{,}96\times0{,}0758/\sqrt{200}\approx0{,}0105$ --- exatamente a margem registrada. O pareamento subtrai a dificuldade de cada tarefa antes de agregar: a diferença média $+0{,}035$ com margem $0{,}011$ não cruza zero.
